\subsection{Hydrography from NHDPlus HR}
\label{sec:methods-nhd}

Watershed delineation and channel geometry for automated SWAT\textsuperscript{+} setup use NHDPlus HR vector products for the conterminous United States \citep{Moore2019NHDPlusHR, Moore2025NHDPlusHR}, together with WBD HU8 and HU12 polygons distributed by vector processing unit (VPU), aligned with HUC4 regions \citep{Jones2022WBD}.
Core layers include flowlines (\texttt{NHDFlowline}), catchments (\texttt{NHDPlusCatchment}), value-added attributes (\texttt{NHDPlusFlowlineVAA}), waterbodies (\texttt{NHDWaterbody}), and WBD HU8/HU12; reach elevations from NHDPlus HR grids support channel drop.
Preprocessing is applied \emph{per VPU} before clipping to user-specified HUC12 lists or analogous extents, so that topology cleaning and hydrologic sequence edits operate on a consistent national graph within each HU4 tile.

NHDPlus HR is not a single-resolution product: its flowlines derive from the National Hydrography Dataset, whose source map scale varies by region (recorded per source in \texttt{NHDSourceCitation.SourceScaleDenominator}). Across our conterminous-U.S.\ archive the finest source scale ranges from 1:1,200 to 1:24,000 and reach density spans roughly an order of magnitude, so SWATGenX records the source resolution and publication vintage per model (Supplementary Table~\ref{tab:nhdplus-provenance}), making delineation density interpretable against the underlying hydrography rather than assumed uniform.

Geometries are projected to a UTM zone from the WBD HU4 centroid (metric lengths and areas), and channel drop is computed from smoothed reach-elevation extrema. The network is then cleaned into a single-valued, tree-like routing graph compatible with the SWAT\textsuperscript{+} subbasin hierarchy \citep{Bieger2017SWATPlus}: divergence-code-2 reaches are removed (their catchment area merged downstream to conserve contributing area), orphan and coastal-artificial flowlines are dropped, the network is partitioned into subbasins satisfying a one-outlet rule, waterbodies above $0.1~\mathrm{km}^2$ are wired to channels under a single-primary-outlet simplification, and a modified-input post-pass merges subbasins below a documented minimum area ($A_{\min}=62.5~\mathrm{km}^2$) into their downstream neighbor where one exists. Full rule-by-rule detail---CRS and geometric attributes, topology cleaning, one-outlet scaffolding, lake--channel linkage, and the $A_{\min}$ minimum-area behavior (including small-domain retention)---is given in Supplementary Section~\ref{sec:supp-nhd-extended}.

\begin{figure}[htbp]
  \centering
  \includegraphics[width=\textwidth]{final/fig-nhd-workflow.png}
  \caption{NHDPlus HR preprocessing stages from VPU-scoped vectors to SWAT\textsuperscript{+} export (Objective~2).}
  \label{fig:nhd-workflow}
\end{figure}

\begin{table}[htbp]
  \centering
  \caption{Core NHDPlus HR preprocessing rules (abridged).}
  \label{tab:nhd-rules-abridged}
  \footnotesize
  \resizebox{\textwidth}{!}{%
    \begin{tabular}{@{}llp{6.8cm}@{}}
\toprule
Rule & Name & Processing action \\
\midrule
NH-00 & VPU / HU4 hydrography unit selection & Parse HU4 listing; download and unpack geodatabase and rasters for the selected processing unit \\
NH-01 & Required NHDPlus HR core layers & Extract required layers to projected working tables following the documented preprocessing s\ldots \\
NH-02 & WBD HU8/HU12 attribution & Spatial join centroids; carry tohuc; multipart dedupe largest area \\
NH-03 & CRS / UTM projection & Select UTM zone from HU4 centroid; project features \\
NH-04 & Divergence code 2 removal & Remove reaches with DivDivergence=2; merge each removed reach catchment into immediate downs\ldots \\
NH-05 & Stream–catchment pairing integrity & Drop flowlines without matching NHDPlusCatchment \\
NH-06 & Coastal / artificial flowline stripping & Filter using implemented coastal and boundary identifier rules \\
NH-07 & Isolated reach policy & Remove from working routed set; optional export to QA folder \\
NH-08 & Reach length and channel drop & Convert km→m; compute drop from smoothed elevations (cm→m scaling) \\
NH-09 & Headwater and terminal flags & Recompute StartFlag and TerminalFl from cleaned links \\
NH-10 & Terminal-reach upstream walk (Subbasin\_level\_1) & Walk upstream from each terminal reach \\
NH-11 & One-outlet partition & Split multi-outlet groups by upstream trace; iterate until $\leq$1 outlet per subbasin withi\ldots \\
NH-12 & Lake / waterbody integration & Area threshold 0.1 km2; map FType to SWAT+ classes; harmonize adjacent lakes; renumber lake/\ldots \\
\bottomrule
\end{tabular}
%
  }
\end{table}

These rules prioritize a single-valued routing graph and consistent, auditable national exports over retaining every national-graph feature; they are not claimed to be hydrologically optimal discretizations, and alternative rule sets would require separate evaluation. Figure~\ref{fig:nhd-workflow} summarizes the stages and Table~\ref{tab:nhd-rules-abridged} lists the core rules (complete catalog: Supplementary Table~S1).

\subsubsection*{Quality assurance checks}
Programmatic checks include one-to-one stream--catchment pairing by \texttt{NHDPlusID}, assignment of each reach to exactly one subbasin, participation of retained lakes with modeled outlets (unresolved cases logged), removal of cyclic constructs on the SWAT\textsuperscript{+} routing, a single outlet per subbasin, and unique exported channel/lake identifiers. Quantitative before/after reach and catchment counts for the evaluation domains appear in Table~\ref{tab:nhd-preprocessing-qa} (\S\ref{sec:results-obj2}); column definitions are given there.

\begin{table}[htbp]
  \centering
  \caption{NHDPlus HR change from original source vectors to exported SWAT\textsuperscript{+} model vectors (Objective~2).}
  \label{tab:nhd-preprocessing-qa}
  \scriptsize
  \setlength{\tabcolsep}{3pt}
  \renewcommand{\arraystretch}{1.05}
  \resizebox{\textwidth}{!}{%
    % Generated by publication/analysis/scripts/emit_tab_nhd_preprocessing_qa_tex.py
% Source: /data/SWATGenXApp/codes/publication/tables/tab-nhd-preprocessing-qa.csv
\begin{tabular}{@{}llrrrrrrrrrr@{}}
\toprule
Size class & Model & Orig. cat. & Orig. reach & Div-2 & Isol. reach & Isol. cat. & Orphan cat. & No-VAA & Final ch. & Final sub. & Final ws. \\
\midrule
S & \texttt{03080102} & 56 & 56 & 2 & 2 & 2 & 4 & 0 & 45 & 3 & 53 \\
M & \texttt{09471300} & 1{,}382 & 1{,}382 & 11 & 0 & 0 & 11 & 0 & 1{,}371 & 12 & 1{,}371 \\
L & \texttt{03100101} & 9{,}739 & 9{,}705 & 194 & 88 & 88 & 282 & 34 & 8{,}181 & 162 & --- \\
\bottomrule
\end{tabular}
%
  }
\end{table}
